\runninghead{Each Text in Its Setting}
\section{Each Text in Its Setting}
\label{sec:each-element}

\subsection{Introit: Jeremias writes to the exiles}
\label{sec:element-introit}

The antiphon comes from a letter. Jeremias sent it from Jerusalem ``to the
residue of the ancients that were carried into captivity, and to the priests,
and to the prophets, and to all the people, whom Nabuchodonosor had carried
away from Jerusalem to Babylon,'' after King Jechonias, the princes of Juda and
the craftsmen had been taken there (Jer 29:1--2). The letter tells the exiles
to settle: ``Build ye
houses, and dwell in them: and plant orchards, and eat the fruit of them''
(v.~5). They are to marry and raise children in Babylon, to ``seek the peace of
the city, to which I have caused you to be carried away captives; and pray to
the Lord for it'' (v.~7), and to give no hearing to the prophets among them who
prophesy falsely in God's name (vv.~8--9). Only then comes the promise: ``When
the seventy years shall begin to be accomplished in Babylon, I will visit
you: and I will perform my good word in your favour, to bring you again to
this place'' (v.~10). The words the Introit sings follow it (vv.~11--14).

The antiphon keeps three things God says he will do: think thoughts of peace,
hear those who call, and bring back the captivity. It opens \latin{Dicit
Dóminus}, ``The Lord saith,'' so that the whole antiphon is God speaking in the
first person, as it is in the letter. It leaves out ``to give you an end and
patience,'' the going and praying, the seeking with the whole heart (v.~13),
and the gathering ``out of all nations.'' St Jerome glosses the omitted ``end
and patience'' as \latin{finem captivitatis, et patientiam laborum
praesentium}, ``the end of the captivity, and patience in the present
labours'' (\work{In Hieremiam}~V, on 29:10--14, PL~24 col.~893); that is his
comment on the words the antiphon does not sing.

The verse, Ps 84:2, opens a psalm ``for the sons of Core'' (84:1) that speaks
first in past tenses: ``Lord, thou hast blessed thy land: thou hast turned away
the captivity of Jacob. Thou hast forgiven the iniquity of thy people: thou
hast covered all their sins'' (vv.~2--3). Two verses later the same psalm
prays, ``Convert us, O God our saviour: and turn off thy anger from us''
(v.~5), and it ends looking forward: ``Surely his salvation is near to them
that fear him \dots\ Justice shall walk before him'' (vv.~10, 14). The Introit
sings only the first of these verses.

\subsection{Collect: thought carried into word and deed}
\label{sec:element-collect}

The Collect has one petition. Its participle, \latin{semper rationabília
meditántes}, ``always meditating on reasonable things,'' describes those who
pray; its main verb, \latin{exsequámur}, asks that they carry out ``what is
pleasing to thee'' (\latin{quæ tibi sunt plácita}) ``both in words and in
deeds'' (\latin{et dictis \dots\ et factis}). The \latin{et \dots\ et} joins
word and deed as a pair. \latin{Rationabília} are things that accord with
reason. The participle describes the thinking as constant (\latin{semper});
the petition is for speech and action. The same Collect is said on the Sixth Sunday
after the Epiphany in its own season.

\subsection{Epistle: a church founded in a few weeks and at once persecuted}
\label{sec:element-epistle}

Paul came to Thessalonica from Philippi, where
he had been ``shamefully treated'' (1~Thess 2:2). He reasoned in the synagogue
for three sabbaths that the Christ had to suffer and rise; some Jews believed,
with ``of those that served God and of the Gentiles a great multitude.'' Then
the Jews who did not believe gathered a mob, set the city in an uproar,
attacked the house of Jason, dragged Jason and some of the brethren before the
rulers, and charged them with saying ``that there is another king, Jesus.''
Jason gave security, and that night the brethren sent Paul and Silas away to
Berea (Acts 17:1--10). Later he sent Timothy from Athens ``to
confirm you and exhort you concerning your faith: That no man should be moved
in these tribulations,'' and wrote this letter when Timothy returned with good
news of their ``faith and charity'' (3:1--7).

The Missal's lesson begins after the greeting of verse~1 and runs to the end
of the chapter. It is one long thanksgiving. Paul remembers three things
before God: \latin{óperis fidei vestræ, et labóris, et caritátis}, ``the work
of your faith and labour and charity,'' and \latin{sustinéntiæ spei}, ``the
enduring of the hope'' (v.~3). In the Latin \latin{labóris} and
\latin{caritátis} stand as two of three coordinated genitives. The Greek text
that St John Chrysostom and Theodoret expound reads the second and third as
one phrase: ``the labour of love.'' Verse~5 says how the gospel came:
\latin{non \dots\ in sermóne tantum, sed et in virtúte}, ``not \dots\ in word
only, but in power also.'' Verse~6
says how they received it: \latin{in tribulatióne multa, cum gáudio Spíritus
Sancti}, ``in much tribulation, with joy of the Holy Ghost.'' Verses~7--8 say
what followed: they became \latin{forma}, a pattern, to the believers of
Macedonia and Achaia, and their faith \latin{profécta est}, ``is gone forth,''
\latin{in omni loco}, ``in every place.'' Verses~9--10 say what is reported of
them: they turned to God from idols ``to serve the living and true God, And to
wait for his Son from heaven \dots\ Jesus, who hath delivered us from the wrath
to come.'' The last verb, \latin{erípuit}, is a perfect: the rescue is spoken
of as done, and the wrath (\latin{ira ventúra}) as still to come.

\subsection{Gradual: thanks remembered in a psalm of complaint}
\label{sec:element-gradual}

Psalm~43 is ``for the sons of Core, to give understanding'' (43:1). It begins
with what the fathers have told: God planted his people in the land, ``For
they got not the possession of the land by their own sword \dots\ But thy
right hand and thy arm'' (vv.~2--4). The speaker takes up the same confidence
for himself, ``I will not trust in my bow: neither shall my sword save me. But
thou hast saved us from them that afflict us \dots\ In God shall we glory all
the day long'' (vv.~7--9). The Gradual is these two verses. The psalm then
turns: ``But now thou hast cast us off, and put us to shame: and thou, O God,
wilt not go out with our armies'' (v.~10). The rest is a complaint of defeat
and scattering. The people protest that they have not forgotten God and ask
why he sleeps: ``Because for thy sake we are killed all the day long: we are
counted as sheep for the slaughter. Arise, why sleepest thou, O Lord?''
(vv.~22--23). St Paul put verse~22 into his letter to the Romans (Rom 8:36);
St Augustine notes that the Apostle quoted it to console a Church suffering
persecution, and says, ``Let us then hear in this Psalm the voice of the
Martyrs'' (\work{En.\ in Ps.}~43.1).

The Gradual keeps the thanks and leaves the complaint. Its \latin{Liberásti
nos} and \latin{confudísti} are past; its verse, \latin{In Deo laudábimur tota
die}, is future; its last words, \latin{in sǽcula}, ``for ever,'' are the
chant's plural where the Vulgate has \latin{in saeculum}.

\subsection{Alleluia: the cry from the depths}
\label{sec:element-alleluia}

Psalm~129 is headed \latin{Canticum graduum}, ``a gradual canticle,'' a title
the chant omits. Its first two verses are a cry: ``Out of the depths I have cried to thee, O Lord: Lord, hear my voice.''
The psalm goes on to say what the depths are and why there is hope in them:
``If thou, O Lord, wilt mark iniquities: Lord, who shall stand it. For with
thee there is merciful forgiveness \dots\ my soul hath hoped in the Lord. From
the morning watch even until night, let Israel hope in the Lord. Because with
the Lord there is mercy: and with him plentiful redemption. And he shall
redeem Israel from all his iniquities'' (vv.~3--8). The Alleluia sings only
the cry, with \latin{oratiónem meam}, ``my prayer,'' where the Vulgate has
\latin{vocem meam}, ``my voice.'' St Hilary's text of the psalm reads
\latin{vocem meam} (\work{Tractatus in Ps.}~129.2--3).

\subsection{Gospel: two parables told to the crowds}
\label{sec:element-gospel}

The parables belong to the discourse of Matthew~13. Jesus sits in a boat and
the crowd stands on the shore (vv.~1--2). He tells the parable of the sower;
the disciples ask why he speaks to the crowds in parables (v.~10), and he
answers, ``to you it is given to know the mysteries of the kingdom of heaven:
but to them it is not given'' (v.~11), quoting Isaias on those who see and do
not perceive (vv.~14--15). He tells the parable of the cockle sown among the wheat, then
these two, and then sends the crowds away and explains the cockle to the
disciples in the house (v.~36). The Missal's lesson begins with the mustard
seed and ends with the evangelist's comment of verses 34--35; it opens with
the liturgical \latin{In illo témpore: Dixit Iesus turbis parábolam hanc}, ``At
that time, Jesus spoke this parable to the multitudes.''

Each parable has something unexpected in it. Mustard is a garden herb,
and the parable says so: when grown it is \latin{maius \dots\ ómnibus
oléribus}, ``greater than all herbs.'' Then it says \latin{et fit arbor}, ``and
becometh a tree,'' big enough for birds to lodge in. Elsewhere in the New
Testament leaven stands for corruption: ``Take heed and beware of the leaven
of the Pharisees'' (Mt 16:6); ``Purge out the old leaven'' (1~Cor 5:7). Here
it is the kingdom of heaven. The woman \latin{abscóndit}, ``hid,'' the leaven
in \latin{tribus satis}, three measures, a large quantity: Jerome notes that
the \latin{satum} was a Palestinian measure holding a \latin{modius} and a half
(\work{In Matthaeum}~II, on 13:33, PL~26 col.~92). The leaven works
\latin{donec fermentátum est totum}, ``until the whole was leavened.''

The evangelist's comment, ``All these things Jesus spoke in parables to the
multitudes: and without parables he did not speak to them'' (v.~34), is
followed by a fulfilment citation. Matthew calls its author ``the prophet''
and quotes Ps 77:2 in a form of his own: the Vulgate's psalm reads ``I will
open my mouth in parables: I will utter propositions from the beginning,''
where Matthew has \latin{eructábo abscóndita a constitutióne mundi}, ``I will
utter things hidden from the foundation of the world.'' The psalm's title is
\latin{Intellectus Asaph}, ``Understanding for Asaph'' (77:1). Jerome reports
that some codices of Matthew named Isaias as the prophet, and others later
dropped the name ``because it was not found in Isaias at all''; he thinks the
evangelist wrote ``by Asaph the prophet'' and that a copyist who did not know
the name changed it to the better-known Isaias (\work{In Matthaeum}~II, on
13:35, PL~26 col.~92). Chrysostom says that the prophet is brought in to
show that Jesus ``is not making a new thing'' (\work{Hom.\ in Mt.}~47.1).

\subsection{Offertory: the cry sung again}
\label{sec:element-offertory}

The Offertory sings the verses of the Alleluia a second time, during the
offering of the gifts, and repeats its first line at the end: \latin{de
profúndis clamávi ad te, Dómine}. The cry from the depths is heard once before
the Gospel and twice at the Offertory.

\subsection{Secret: four verbs of one petition}
\label{sec:element-secret}

\latin{Hæc nos oblátio, Deus, mundet, quǽsumus, et rénovet, gubérnet et
prótegat}: ``May this oblation, O God, cleanse, renew, govern, and protect
us.'' The subject is the offering (\latin{oblátio}) on the altar; the object
is ``us'' (\latin{nos}), those who offer it. The four verbs come in two pairs,
and the Latin joins each pair with \latin{et}: cleanse and renew, govern and
protect. The first pair asks for a change in those who offer; the second for
their keeping afterwards.

\subsection{Communion: the promise after the withered fig tree}
\label{sec:element-communion}

In Mark the saying stands on the road from Bethany to Jerusalem in the last
days before the Passion. The day before, Jesus had cursed a fig tree that had
leaves and no fruit, and had driven the sellers out of the temple, saying that
it should be ``the house of prayer to all nations'' (Mk 11:12--17). In the
morning Peter sees the fig tree ``dried up from the roots,'' and Jesus answers:
``Have the faith of God. Amen I say to you that whosoever shall say to this
mountain, Be thou removed and be cast into the sea, and shall not stagger in
his heart, but believe that whatsoever he saith shall be done; it shall be done
unto him. Therefore I say unto you, all things, whatsoever you ask when ye
pray, believe that you shall receive: and they shall come unto you. And when
you shall stand to pray, forgive, if you have aught against any man: that your
Father also, who is in heaven, may forgive you your sins'' (Mk 11:20--25). The
antiphon gives verse~24 between the opening words \latin{Amen dico vobis} and
the close \latin{et fiet vobis}. Read alone, it promises whatever is asked to
whoever believes. In Mark the promise follows the saying on mountain-moving
faith and is followed by the command to forgive.

\subsection{Postcommunion: fed, and still hungry}
\label{sec:element-postcommunion}

\latin{Cæléstibus, Dómine, pasti delíciis}: the prayer opens with a past
participle, ``having been fed with heavenly delights.'' Its petition is
\latin{ut semper éadem, per quæ veráciter vívimus, appetámus}, ``that we may
always long for the same things by which we truly live.'' \latin{Éadem}, ``the
same things,'' refers back to the delights just received; \latin{veráciter
vívimus}, ``we truly live,'' says that these are the things of life; and
\latin{appetámus} asks for desire. The Lasance renders
it ``that we may ever hunger after those things.''
